% Bare tabular. Wrap in \begin{table}...\caption{...}\label{tab:inference_sens_high}\end{table} at the call site.
\begin{tabular}{lccccccc}\hline\hline
Inference scheme & h=0 & h=5 & h=20 & h=30 & h=40 & h=50 & h=60 \\\hline
Coefficient (scaled) & -0.323 & -0.224 & -0.467 & -4.139 & -5.898 & -7.153 & -7.711 \\
\hline
Two-way cluster (firm $\times$ event) [preferred] & ( 0.529) & ( 0.872) & ( 1.179) & ( 1.747) & ( 2.067) & ( 2.410) & ( 2.644) \\
Firm cluster & ( 0.325) & ( 0.670) & ( 1.167) & ( 1.821) & ( 2.061) & ( 2.170) & ( 2.283) \\
Event cluster & ( 0.531) & ( 0.910) & ( 1.167) & ( 1.493) & ( 1.810) & ( 2.254) & ( 2.545) \\
Heteroskedasticity-robust & ( 0.326) & ( 0.718) & ( 1.153) & ( 1.571) & ( 1.794) & ( 1.987) & ( 2.160) \\
\hline
Wild-cluster bootstrap by event: 95\% CI &   &   &   & [-7.24, -1.22] & [-9.43, -2.36] & [-11.56, -2.75] & [-12.68, -2.74] \\
\hline\hline
\multicolumn{8}{l}{\footnotesize Point estimate $\hat\beta \times 25 \times 100$ (in percentage points) is constant across rows by construction; only the variance estimator differs.}
\multicolumn{8}{l}{\footnotesize Two-way clustering by firm and event is the preferred specification. Firm- and event-only clustering are nested in the two-way scheme. Heteroskedasticity-robust is reported as a lower bound on the SE.}
\multicolumn{8}{l}{\footnotesize Wild-cluster bootstrap by event ($9{,}999$ replications, percentile-$t$ CI) addresses the small number of event clusters ($n = 159$). Reported at the four long horizons where the headline differential is statistically distinguishable from zero in the preferred specification.}
\end{tabular}
